\section{Front wheeler: method and cost model}
\label{sec:framework}

\paragraph{Notation.} Let $f(x) = \mathrm{head}(\mathrm{trunk}(x;\tht);\thh)$
be a depth-$(D{+}1)$ network, where $\thh$ collects the \emph{front}---the
last hidden layer (or last transformer block) and the output head---and
$\tht$ the trunk below. Training minimizes a loss $L(\thh,\tht)$.

\paragraph{Backward cost.} We count backward multiply--accumulates per
weight matrix as ${\sim}2\,ab$ for a matrix of shape $a\times b$ (one GEMM
for the activation gradient, one for the weight gradient). A \emph{full} step
back-propagates through all matrices. A \emph{front} step detaches the trunk
output and back-propagates through the front only---a factor ${\sim}(D{+}1)/2$
cheaper for uniform widths, and the trunk's no-grad forward additionally
skips activation storage. A \emph{trunk-only} step (the \bcd{} ablation)
must still compute activation gradients \emph{through} the head to reach the
trunk: it saves only the head's weight-gradient GEMMs and costs nearly as
much as a full step. This asymmetry is the crux: depth truncation is
available only to the top block, so partitioned block-coordinate schemes
cannot convert block restriction into compute savings, whereas the nested
scheme (front vs.\ everything) can. Block-coordinate theory assumes uniform
per-block costs \citep{nutini2015coordinate}; the cost asymmetry \emph{is}
the scheduling problem.

\paragraph{Terminology and comparison protocol.} Schedules are written
with their parameters: \textsc{periodic}-$N$ (one full step every $N$),
\lpft{}($s$) (front until fraction $s$ of training, then full),
\textsc{ratchet}($r$) (all-but-front frozen by fraction $r$ of training),
\fw{}($\tau$, $\kappa_{\max}$) (trigger threshold and cooldown cap), and
\bcd{} (as \fw{} with trunk-only bursts). Cost is reported on two axes:
\emph{deep-step budget} (fraction of steps at full depth; meaningful only
for schedules whose steps are binary front/full) and \emph{backward
compute} (total backward MACs, as a fraction of \textsc{full}'s; defined
for every schedule). All cross-schedule comparisons in this paper are made
at matched backward compute; deep-step budget is reported for intuition.

\paragraph{The schedule as a controller.} Training chooses per step between
a front step and a full step; the budget is the fraction of full-depth
steps. The greedy, cost-normalized rule would take front steps while
\begin{equation}
\frac{\|P\nabla L\|^2}{c_{\mathrm{front}}} \;\ge\;
\frac{\|\nabla L\|^2}{c_{\mathrm{full}}},
\label{eq:greedy}
\end{equation}
where $P$ projects onto the front coordinates---a cost-aware
Gauss--Southwell rule over \emph{nested} blocks. Neither side of
\eqref{eq:greedy} is free to observe: the left is estimated by realized
front-phase loss improvement, while the right requires paying
$c_{\mathrm{full}}$. \fw{} (Algorithm~\ref{alg:fw}) therefore treats each
burst as both a progress step and a \emph{measurement}: an unproductive
burst reveals that the plateau is global, and exponential backoff raises the
evidence threshold for the next alarm---the ``boy who cried wolf''
correction. We call the cap $\kappa_{\max}$ on this backoff the
\emph{governor}; Section~\ref{sec:exp3} shows it doubles as the budget
dial. Bursts are \emph{consecutive} full steps by design; the
experiments show that consecutiveness is not an implementation convenience
but a stability requirement (Sections~\ref{sec:exp2}, \ref{sec:exp3}).
